
\section{Applications}
\label{sec:6.8}

\subsection{Prime Cycle}
\label{sec:6.8.1}

If $c$, the length of a single limiting cycle for $R$, is prime, 
then the Fourier condition means that that the first row of $\Sigma(\pi,\R)'$ is not constant. 
Using BM-F, we know there is an aperiodic vertex. 
If not already synchronized, recoloring there guarantees that the probabilities on the cycle $C$ are nonconstant. 
Hence there is a kernel which is rank one.

\subsection{Friedman}
\label{sec:6.8.2}
% BEGIN SUBJECT INDEX: sec:6.8.2
\index{Friedman weight}
% END SUBJECT INDEX


Let $W$ denote the \textsl{Friedman weight\/}. I.e., rationalize $\pi$ to $\tilde\pi$ which has integer entries
with gcd equal to 1. These are the \textsl{weights\/}. The Friedman weight is the sum of those weights.
From BM-F it immediately follows that $r$, the rank of the kernel, divides $W$.
If $c$ and $W$ are relatively prime, then since $r$ divides both $c$ and $W$, the kernel $\cK$ has rank one.
